% Chapter — Cannonball orbital drag (FR-9). Follows the mandatory FR-29
% six-section template: purpose; assumptions and domain of validity;
% derivation; discretization and implementation; validation evidence;
% references. The equation labels eq:drag:accel and eq:drag:vrel are
% echoed verbatim in comments in cpp/src/models/drag.cpp (FR-29
% equation-label-to-code traceability); renaming them breaks that trace.
\chapter{Cannonball Orbital Drag}
\label{ch:drag}

\section{Purpose}
\label{sec:drag:purpose}

This chapter documents the cannonball atmospheric-drag acceleration
(FR-9): the orbital-regime drag model that consumes a density from one of
the atmosphere models (Chapters~\ref{ch:ussa76},
\ref{ch:harrispriester}, and~\ref{ch:marsatmosphere}) and a ballistic
parameter $C_D A / m$, mirroring the cannonball solar-radiation-pressure
model's $C_R A / m$. FR-9 splits the aerodynamic domain deliberately: the
Mach-table aerodynamic database (Phase~4) is valid for continuum ascent
flight only, while this cannonball model serves the orbital free-molecular
regime; the two never apply simultaneously, and a vehicle selects the
cannonball model per configuration. The model is implemented in
\texttt{cpp/src/models/drag.cpp} with its interface in
\texttt{cpp/include/star/models/drag.hpp}.

\section{Assumptions and domain of validity}
\label{sec:drag:domain}

Assumptions:
\begin{enumerate}
  \item \textbf{Cannonball geometry.} The drag force is
    attitude-independent: a single constant product $C_D A / m$
    characterizes the vehicle, as for a sphere. For non-spherical
    vehicles this is the standard mean-area approximation; the error is
    absorbed into the configured $C_D A / m$.
  \item \textbf{Co-rotating atmosphere (the FR-8 air-relative-velocity
    rule).} All aerodynamic quantities use the velocity of the vehicle
    relative to an atmosphere rotating rigidly with the planet,
    equation~\eqref{eq:drag:vrel}. Winds are neglected; near-surface wind
    speeds (tens of \unit{\metre\per\second}) are two orders below
    orbital air-relative speeds, and thermospheric super-rotation is far
    inside the density model's own uncertainty.
  \item \textbf{Free-molecular drag coefficient.} In the orbital regime
    the flow is free-molecular (mean free path far exceeds vehicle
    dimensions; cf.\ the USSA76 mean-free-path tables above
    \qty{100}{\kilo\metre}~\cite{ussa1976}). The commonly adopted default
    for compact satellites is $C_D = 2.2$, per the drag discussions of
    Montenbruck and Gill~\cite[Sect.~3.5]{montenbruck2000} and
    Vallado~\cite[Sect.~8.6.2]{vallado2013}; the value enters only
    through the configured $C_D A / m$ (the default itself is applied at
    the vehicle-configuration level in Phase~4, not inside this
    function). Physical free-molecular values for non-compact shapes
    range roughly 2.0--3.0; treat absolute decay rates accordingly.
  \item \textbf{Frame-agnostic kinematics.} The model is a pure function
    of $(\rho,\, C_D A/m,\, \vv{v}_{\mathrm{rel}})$; the EOM layer owns
    frames and supplies $\vv{v}_{\mathrm{rel}}$ resolved in the
    integration frame, so the returned acceleration is in that same
    frame.
\end{enumerate}

Domain of validity: orbital free-molecular flight around any body whose
atmosphere model supplies $\rho$; non-negative $\rho$ and $C_D A / m$.
Out-of-domain behavior, matching the code exactly: negative $\rho$ or
negative $C_D A / m$, and non-finite inputs, throw
\texttt{std::domain\_error}; a zero relative velocity returns the exact
zero vector (no normalization singularity exists because the formulation
never divides by $\lVert \vv{v}_{\mathrm{rel}} \rVert$). Continuum ascent
aerodynamics are explicitly outside this model's domain --- they are the
Phase~4 Mach-table database's regime (FR-9 domain split).

\section{Derivation}
\label{sec:drag:derivation}

A surface element moving at speed
$v = \lVert \vv{v}_{\mathrm{rel}} \rVert$ through gas of density $\rho$
intercepts momentum flux $\rho v^2$ per unit area. Writing the net
aerodynamic force on a body of reference area $A$ as this flux times the
dimensionless drag coefficient $C_D/2$ (the factor 2 is the classical
convention that normalizes the force by the dynamic pressure
$\tfrac{1}{2}\rho v^2$), directed opposite the air-relative velocity
\cite[Sect.~3.5]{montenbruck2000}, \cite[Sect.~8.6.2]{vallado2013}:
\begin{equation}
  \vv{F} \;=\; -\tfrac{1}{2}\,\rho\,C_D A\,
    v\,\vv{v}_{\mathrm{rel}},
\end{equation}
where $v\,\vv{v}_{\mathrm{rel}} = v^2\,\hat{\vv{v}}_{\mathrm{rel}}$
combines the magnitude and direction without a normalization. Dividing by
the vehicle mass gives the model's governing equation,
\begin{equation}
  \boxed{\;
  \vv{a}_{\mathrm{drag}} \;=\;
    -\tfrac{1}{2}\,\rho\,\frac{C_D A}{m}\,
    \lVert \vv{v}_{\mathrm{rel}} \rVert\, \vv{v}_{\mathrm{rel}}\; }
  \label{eq:drag:accel}
\end{equation}
with the air-relative velocity of a co-rotating atmosphere (FR-8)
\begin{equation}
  \vv{v}_{\mathrm{rel}} \;=\; \vv{v} \;-\;
    \vv{\omega}_{p} \times \vv{r},
  \label{eq:drag:vrel}
\end{equation}
where $\vv{r}, \vv{v}$ are the vehicle's planet-centered inertial position
and velocity and $\vv{\omega}_{p}$ is the planet's angular-velocity vector
(Earth: $\omega_\oplus = \qty{7.292115e-5}{\radian\per\second}$, IERS
Conventions~\cite{iers2010}; Mars: $\omega_{\mathrm{Mars}} =
\qty{7.088218e-5}{\radian\per\second}$, from the IAU 2015 rotation rate of
Archinal et al.~\cite{archinal2018} --- both defined in
\texttt{star/constants.hpp}). Equation~\eqref{eq:drag:vrel} is evaluated
by the EOM layer, which owns frames; this module receives
$\vv{v}_{\mathrm{rel}}$ ready-made, and the same air-relative velocity
rule governs the Phase~4 ascent aerodynamics (Mach number and dynamic
pressure), so the two drag regimes share one kinematic convention.

\section{Discretization and implementation}
\label{sec:drag:implementation}

\begin{enumerate}
  \item \textbf{Module and traceability.} The model lives in
    \texttt{cpp/src/models/drag.cpp}; the source comments echo the labels
    \texttt{eq:drag:accel} and \texttt{eq:drag:vrel} verbatim.
  \item \textbf{Single expression, no singularity.} The implementation is
    the literal product
    $-\tfrac{1}{2}\rho\,(C_D A/m)\,
    \lVert\vv{v}_{\mathrm{rel}}\rVert\,\vv{v}_{\mathrm{rel}}$ using
    Eigen's \texttt{norm()}; at zero velocity the product is exactly zero
    with no special case.
  \item \textbf{Determinism.} Three multiplies and a norm per component
    in fixed order; the only libm call is the square root inside the
    norm. Same-platform bit identity is trivial; cross-platform spread
    is one ulp of \texttt{sqrt}.
\end{enumerate}

\section{Validation evidence}
\label{sec:drag:validation}

Table~\ref{tab:drag:validation} lists the tests that constitute this
chapter's validation evidence. Each test identifier is machine-checked to
exist in the test tree by \texttt{scripts/lint\_chapters.py}; golden
provenance and tolerances are recorded in
\texttt{tests/golden/atmosphere/manifest.toml}.

\begin{table}[htbp]
  \centering
  \caption{Validation evidence for the cannonball drag model (doctest,
    \texttt{cpp/tests/test\_drag.cpp}).}
  \label{tab:drag:validation}
  \begin{tabular}{lp{8.6cm}}
    \toprule
    Test ID & Acceptance basis \\
    \midrule
    \testid{DRAG-CANNONBALL-GOLDEN} &
      Five committed states (axis-aligned, generic, dense-ascent,
      bulge-state, zero-velocity) match the independent 50-digit mpmath
      reference component-wise to $10^{-14}$ relative; the zero-velocity
      case returns the exact zero vector. \\
    \testid{DRAG-PROPERTIES} &
      FR-22 layer-2 invariants of equation~\eqref{eq:drag:accel}:
      acceleration is anti-parallel to $\vv{v}_{\mathrm{rel}}$
      (cross product exactly zero, dot product negative), linear in
      $\rho$ and in $C_D A/m$, and quadratic under velocity scaling
      ($\vv{a}(2\vv{v}) = 4\,\vv{a}(\vv{v})$ to double rounding);
      domain errors per Section~\ref{sec:drag:domain}. \\
    \bottomrule
  \end{tabular}
\end{table}

\section{References}
\label{sec:drag:references}

This chapter relies on Montenbruck and Gill~\cite{montenbruck2000},
Sect.~3.5, and Vallado~\cite{vallado2013}, Sect.~8.6.2, for the drag
formulation and the conventional $C_D = 2.2$ default; on the IERS
Conventions (2010)~\cite{iers2010} for the Earth rotation rate; and on
Archinal et al.~\cite{archinal2018} for the Mars rotation rate. Full
bibliographic details appear in the bibliography.
